\documentclass[tikz, border=8pt]{standalone}
\usepackage[utf8]{inputenc}
\usepackage[T1]{fontenc}
\usepackage[french]{babel}
\usetikzlibrary{positioning, arrows.meta, shapes.geometric, calc, shapes.multipart}

\definecolor{dbblue}{RGB}{10, 22, 40}

\tikzset{
  entity/.style={
    draw=dbblue, line width=0.9pt,
    rectangle split, rectangle split parts=2,
    rectangle split part fill={dbblue, white},
    rectangle split part align={center, left},
    text width=5.7cm, font=\footnotesize,
    inner sep=4pt,
  },
  rel/.style={draw=black!60, line width=0.9pt, -{Stealth[length=6pt]}},
  relD/.style={draw=black!60, line width=0.8pt, dash pattern=on 3pt off 2pt},
  lbl/.style={font=\scriptsize\itshape, fill=white, inner sep=1pt},
}
\newcommand{\eh}[1]{\textcolor{white}{\normalfont\bfseries\small #1}}
\newcommand{\ea}[1]{\textcolor{black}{\small #1}}
\newcommand{\id}[1]{\underline{#1}}
\newcommand{\fk}[1]{#1}

\begin{document}
\begin{tikzpicture}[>=Stealth]

\node[entity] (transp) at (0,9) {\eh{TRANSPORTEURS}
\nodepart{two}\ea{
\id{id} : Entier (PK)\\
nom : Chaine\\
contact : Chaine}};

\node[entity] (tournees) at (0,5.0) {\eh{TOURNEES}
\nodepart{two}\ea{
\id{id} : Entier (PK)\\
\fk{transporteur\_id} : Entier (FK)\\
date : Date\\
vehicule\_id : Chaine\\
chauffeur : Chaine}};

\node[entity] (livraisons) at (0,0.2) {\eh{LIVRAISONS}
\nodepart{two}\ea{
\id{id} : Entier (PK)\\
\fk{tournee\_id} : Entier (FK)\\
tracking\_number : Chaine\\
adresse\_livraison : Chaine\\
heure\_estimee : Chaine\\
heure\_reelle : Chaine}};

\node[entity] (expeditions) at (11.2,0.2) {\eh{EXPEDITIONS}
\nodepart{two}\ea{
\id{id} : Entier (PK)\\
commande\_id : Entier (FK externe)\\
tracking\_number : Chaine\\
transporteur : Chaine\\
date\_expedition : Date\\
date\_livraison\_prevue : Date\\
date\_livraison\_reelle : Date\\
statut : Chaine}};

\node[entity] (cc) at (8.0,9) {\eh{COMMANDES\_CLIENTS}
\nodepart{two}\ea{
\id{id} : Entier (PK)\\
client : Chaine\\
commande\_id : Chaine (metier)\\
date\_commande : Date\\
entrepot : Chaine}};

\node[entity] (lcc) at (8.0,5.0) {\eh{LIGNES\_COMMANDE\_\\CLIENTS}
\nodepart{two}\ea{
\id{id} : Entier (PK)\\
\fk{commande\_client\_id} : Entier (FK)\\
sku : Chaine (metier)\\
quantite : Entier}};

\node[entity] (produits) at (18.0,9) {\eh{PRODUITS}
\nodepart{two}\ea{
\id{id} : Entier (PK)\\
sku : Chaine\\
libelle : Chaine\\
categorie : Chaine\\
poids\_kg : Decimal\\
temperature\_dirigee : Booleen}};

% --- FK formelles (fleche vers la PK referencee) ---
\draw[rel] (tournees.north) -- (transp.south)
  node[lbl, pos=0.5, right=2pt] {FK transporteur\_id -> id};

\draw[rel] (livraisons.north) -- (tournees.south)
  node[lbl, pos=0.5, right=2pt] {FK tournee\_id -> id};

\draw[rel] (lcc.north) -- (cc.south)
  node[lbl, pos=0.5, right=2pt] {FK commande\_client\_id -> id};

% --- Rapprochements par cle metier, sans contrainte FK ---
\draw[relD] (livraisons) -- (expeditions)
  node[lbl, pos=0.5, above=2pt] {tracking\_number (cle metier, pas de FK)};

\draw[relD] (produits.south west) -- (lcc.north east)
  node[lbl, pos=0.78, above right=1pt] {sku (cle metier, pas de FK)};

\node[font=\footnotesize, align=left, below=1.3cm of livraisons, xshift=6.0cm] (legend) {
  \tikz[baseline=-0.5ex]{\draw[rel] (0,0) -- (0.9,0);} contrainte FK (SGBD) \qquad
  \tikz[baseline=-0.5ex]{\draw[relD] (0,0) -- (0.7,0);} rapprochement par cle metier (pas de FK)
};

\node[font=\footnotesize\itshape, align=left, below=0.35cm of legend] {
  Memes 7 tables que le MCD en regard -- types generiques (MLD), encore independants de PostgreSQL (MPD reel : schema\_base\_travail.sql)
};

\end{tikzpicture}
\end{document}
